

\section{ORGANIZAÇÃO DA DISSERTAÇÃO}

Esta dissertação está estruturada em cinco capítulos principais, organizados de forma a conduzir o leitor desde os fundamentos teóricos até o desenvolvimento das contribuições metodológicas propostas, culminando em uma discussão consolidada dos resultados obtidos.

\textcolor{red}{Rever a estrutura. Minha sugestao eh ter um capitulo de Avaliação de Agrupamente; e um Capitulo de TRI, onde vc incluiria o Beta4. Fica mais organizado por assunto.}

\textcolor{blue}{Correção: estrutura reorganizada por assunto. O Capítulo \ref{chap:n2_agrupamentos} trata da avaliação de agrupamentos; o Capítulo \ref{chap:n3_irt} trata da TRI e inclui o $\beta^{4}$-\gls{IRT} como seção; a comparação empírica $\beta^{3}$ vs. $\beta^{4}$ foi movida para o Apêndice \ref{ap:beta3beta4}.}

O Capítulo \ref{chap:n2_agrupamentos} discute o problema de validação de agrupamentos, detalhando índices de validação internos, tais como \textit{Calinski-Harabasz}, \textit{Silhouette} e \textit{Davies-Bouldin}, bem como índices externos, como \textit{Mutual Information}, \textit{V-measure} e \textit{Adjusted Rand Index}.

O Capítulo \ref{chap:n3_irt} apresenta a fundamentação teórica que sustenta o desenvolvimento das propostas. Inicialmente, é introduzida a \gls{TRI}, abordando seus principais conceitos, parâmetros e modelos clássicos. Em seguida, apresenta-se o modelo $\beta^{3}$-\gls{IRT}, que serve como base para a extensão proposta na seção seguinte, na qual é apresentado o modelo $\beta^{4}$-\gls{IRT}, com a reformulação do parâmetro de discriminação em componentes de sinal e magnitude e a comparação teórica e empírica com o modelo $\beta^{3}$-\gls{IRT}, cujos protocolos experimentais e resultados detalhados são apresentados no Apêndice \ref{ap:beta3beta4}.

O Capítulo \ref{chap:n4_claire} apresenta o método \gls{CLAIRE}, um \textit{framework} para avaliação de agrupamentos baseado em \gls{TRI} e na concordância entre modelos. Inicialmente, é definida a formulação do modelo e discutida a relação entre concordância e corretude. Em seguida, são apresentados os experimentos realizados, incluindo o ranqueamento de modelos, a análise da dificuldade e discriminação das instâncias, a correlação entre a habilidade estimada e métricas tradicionais de validação, o efeito da inserção de partições aleatórias no conjunto de modelos e a análise da estabilidade das habilidades estimadas.

Por fim, o Capítulo \ref{chap:n5_conclusao} apresenta as conclusões gerais da dissertação. São sintetizadas as principais contribuições teóricas e metodológicas do trabalho, discutidas suas limitações e delineadas possíveis direções para trabalhos futuros.